The design requires that we do not sanitize the journey into a success
story. This is the honest record of everything that went wrong, was
surprising, or was a failed attempt --- because each of these taught us
something, and several are as valuable as any success.

Each entry: what happened, what we learned, and what we did about it.

\begin{center}\rule{0.5\linewidth}{0.5pt}\end{center}

\section{The evaluator reported the model at 0.000 --- but the model was
fine}\label{the-evaluator-reported-the-model-at-0.000-but-the-model-was-fine}

\textbf{What happened.} On the very first Stage-1 evaluation, the
trained model scored a reported \textbf{0.000} (syntax failure on every
problem). This looked like a total failure of the approach.

\textbf{What we found.} The model was actually producing correct
programs --- but our generation output appended null-byte padding, and
Python's parser rejects null bytes. So the \textbf{judge} was scoring
valid programs as syntax errors.

\textbf{What we learned.} The most important lesson of the book,
demonstrated on ourselves: \emph{never trust a number from a system we
have not audited.} Fixing the judge revealed the model's true 0.897. If
we'd taken 0.000 at face value, we'd have wrongly concluded the approach
failed.

\textbf{Action.} Fixed \texttt{generate\_code} to truncate at the first
null byte; re-ran; got 0.897.

\begin{center}\rule{0.5\linewidth}{0.5pt}\end{center}

\section{The pilot wording leaked the answer (100\%
predictability)}\label{the-pilot-wording-leaked-the-answer-100-predictability}

\textbf{What happened.} The leakage audit trained a trivial TF-IDF
classifier to predict the operation from the natural-language prompt. It
succeeded \textbf{100\%} of the time.

\textbf{What we learned.} Our first renderer's wording was so tied to
the operation that a model ``succeeding'' might just be exploiting the
phrasing, not learning the computation. This is \emph{why} the
surface-generalization tests matter.

\textbf{Action.} Recorded as finding-001; noted the pilot is fine for
Stage 1 but must add wording-family holdouts before any surface test.

\begin{center}\rule{0.5\linewidth}{0.5pt}\end{center}

\section{A bigger model did NOT fix
composition}\label{a-bigger-model-did-not-fix-composition}

\textbf{What happened.} We scaled the model across a \textasciitilde33×
range (0.88M → 29.3M) expecting the composition boundary to move. It did
not: composition stayed at \textbf{0.000} at all four sizes (0.88M,
6.5M, 10.2M, 29.3M).

\textbf{What we learned.} At \textasciitilde1--30M parameters,
compositional generalization failure is \textbf{not a capacity ceiling}.
It is a structural/learning limit that scaling alone does not resolve.
This is a genuine negative result that redirects the investigation.

\textbf{Action.} Recorded in
\texttt{docs/experiments/scaling\_composition.md}.

\begin{center}\rule{0.5\linewidth}{0.5pt}\end{center}

\section{Sequential training caused catastrophic forgetting (A →
0.000)}\label{sequential-training-caused-catastrophic-forgetting-a-0.000}

\textbf{What happened.} Teaching the model lists after arithmetic,
sequentially, completely destroyed the arithmetic (A-retention
\textbf{0.000}).

\textbf{What we learned.} Catastrophic forgetting is real and severe
here --- but it is a \textbf{procedure} problem, not a fate. Replay
(interleaving A and B) retained A perfectly (1.000) with zero cost to
learning B.

\textbf{Action.} Recorded; the chapter states this clearly --- the
protocol, not capacity, determines the outcome.

\begin{center}\rule{0.5\linewidth}{0.5pt}\end{center}

\section{The pilot task was ``too easy'' (retrieval solved
it)}\label{the-pilot-task-was-too-easy-retrieval-solved-it}

\textbf{What happened.} The retrieval baseline scored \textbf{1.000} on
the pilot test set, and the model scored 0.897. So a non-learning lookup
did better than the model.

\textbf{What we learned.} The pilot task, as constructed, is
retrieval-solvable --- so it does not yet test whether the model
\emph{generalizes}, only whether it matches surface form. This is
exactly the ``task too easy'' outcome the plan anticipated.

\textbf{Action.} We did not hide it; we used it to justify harder splits
(surface, then composition), which produced the real findings above.

\begin{center}\rule{0.5\linewidth}{0.5pt}\end{center}

\section{The human-written test set:
0/16}\label{the-human-written-test-set-016}

\textbf{What happened.} The trained model solved \textbf{0/16}
naturally-phrased human prompts --- even simple ``add two numbers''
phrased as a person would.

\textbf{What we learned.} This is the external, real-world corroboration
of the surface finding. The model learned a \emph{template}, not the
\emph{task}. The moment a real user phrases a request differently, it
fails.

\textbf{Action.} Recorded as an external sanity check; it sharpens the
book's honest conclusion.

\begin{center}\rule{0.5\linewidth}{0.5pt}\end{center}

\section{The free-model proofread returned degraded
output}\label{the-free-model-proofread-returned-degraded-output}

\textbf{What happened.} We dispatched the 85K-char draft to two free
models for a proofread. One leaked its chain-of-thought and broke into
noise; the other timed out. No usable independent review was collected.

\textbf{What we learned.} Free endpoints are flaky for long documents,
and we should not fabricate fixes from unusable output. We were honest
about this rather than pretend a review happened.

\textbf{Action.} Noted honestly; the book's line-edit stands on its own
merit, not a fabricated review.

\begin{center}\rule{0.5\linewidth}{0.5pt}\end{center}

\section{Transient SIGILL crashes during
training}\label{transient-sigill-crashes-during-training}

\textbf{What happened.} Long CPU training runs occasionally crashed with
signal 11 (SIGILL) on the AVX-only Ivy Bridge CPU --- intermittently,
one run succeeds and the next dies.

\textbf{What we learned.} Hardware quirks are real and part of the brew.
The fix is \texttt{TORCH\_DISABLE\_AVX2=1} plus a retry loop (and
\texttt{OMP\_NUM\_THREADS=1} for larger models).

\textbf{Action.} Wrapped all long runs in a retry loop; documented the
workaround.

\begin{center}\rule{0.5\linewidth}{0.5pt}\end{center}

\section{Parameter-math and naming mistakes (caught by
testing)}\label{parameter-math-and-naming-mistakes-caught-by-testing}

\textbf{What happened.} Early writing had a parameter-count
mis-derivation (didn't sum to the true 884K/884,224) and the human-test
evaluator initially referenced a non-existent key.

\textbf{What we learned.} Always verify arithmetic against the real
build, and let the ``test the tester'' discipline catch even our own
book's numbers.

\textbf{Action.} Recomputed from the actual build; fixed the evaluator.

\begin{center}\rule{0.5\linewidth}{0.5pt}\end{center}

\emph{The point of this ledger is the opposite of marketing. Every entry
above, in its own way, made the final understanding sharper --- and
several of them --- the broken evaluator, the wording leak, the
task-too-easy result, the 0/16 human test --- are among the most
valuable findings in the book precisely because they forced us to be
honest.}
